%% kapa: exam=1 round=19 year=2008 period=1 form=A subject=accounting subject_ko=회계학 seq=18 number=98
%% source: Q-Net 감정평가사 기출문제 (cst00309) artlSeq=1001962 file=extracted/1008/1001962/2008년 제19회 감정평가사 1차시험 기출문제/1교시.hwp
%% conversion: kordoc 4.24.51 -> tools/gen1.py (auto); review=proofread-sonnet; flags=none
\begin{question}{98}
유가증권의 양도로 유가증권 보유자가 유가증권의 통제를 상실한 때에는 그 유가증권을 대차대조표에서 제거한다. 다음 중 유가증권의 통제를 상실한 경우에 해당하는 것을 모두 고른 것은?
\begin{framedtext}
ㄱ. 유가증권의 경제적 효익을 획득할 수 있는 권리를 처분한 때 \par
ㄴ. 유가증권의 경제적 효익을 획득할 수 있는 권리가 만료된 때 \par
ㄷ. 유가증권의 경제적 효익을 획득할 수 있는 권리를 전부 실현한 때
\end{framedtext}
\begin{choices}
  \item ㄱ
  \item ㄱ, ㄴ
  \item ㄱ, ㄷ
  \item ㄴ, ㄷ
  \item ㄱ, ㄴ, ㄷ
\end{choices}
\end{question}
\answer{5}
